\documentclass[11pt,oneside]{report}
% ======================================================================
%  THE TRIADIC MEASUREMENT SUBSTRATE --- COMPLETE CORPUS (Volumes 1 & 2)
%  Aaron Switzer, 2026
%  Single-file build. Compile: pdflatex x3 (third pass settles the TOC).
% ======================================================================
% ======================================================================
%  TMS Volume 1 -- shared preamble
%  Compiles with pdfLaTeX both locally and on Overleaf (no exotic fonts).
% ======================================================================
\usepackage[T1]{fontenc}
\usepackage[utf8]{inputenc}
\usepackage{amsmath}
\usepackage{amssymb}
\usepackage{mathptmx}            % Times text + math (widely available)
\usepackage[scaled=0.90]{helvet} % sans for headings
\usepackage{courier}

\usepackage[letterpaper,margin=1in]{geometry}
\usepackage{amsthm}
\usepackage{mathtools}
\usepackage{bm}
\usepackage{microtype}
\usepackage{enumitem}
\usepackage{booktabs}
\usepackage{array}
\usepackage{longtable}
\usepackage{caption}
\usepackage{xcolor}
\usepackage{titlesec}
\usepackage{fancyhdr}
\usepackage[hidelinks]{hyperref}
\usepackage{tocbibind}

% ---- spacing -----------------------------------------------------------
\setlength{\parindent}{1.2em}
\setlength{\parskip}{0.35em}
\linespread{1.04}
\setlist{leftmargin=1.6em,topsep=0.25em,itemsep=0.15em,parsep=0pt}

% ---- theorem environments ---------------------------------------------
\newtheorem{theorem}{Theorem}[section]
\newtheorem{lemma}[theorem]{Lemma}
\newtheorem{corollary}[theorem]{Corollary}
\newtheorem{proposition}[theorem]{Proposition}

\theoremstyle{definition}
\newtheorem{definition}[theorem]{Definition}
\newtheorem{axiom}{Axiom}
\newtheorem{claim}[theorem]{Claim}
\newtheorem{principle}[theorem]{Principle}
\newtheorem{conjecture}[theorem]{Conjecture}
\newtheorem{prediction}[theorem]{Prediction}
\newtheorem{property}[theorem]{Property}
\newtheorem{postulate}[theorem]{Postulate}

\theoremstyle{remark}
\newtheorem{remark}[theorem]{Remark}
\newtheorem*{remark*}{Remark}

% ---- section heading look (unnumbered; author supplies the numbers) ----
\titleformat{\section}[hang]{\normalfont\large\bfseries\sffamily}{}{0pt}{}
\titleformat{\subsection}[hang]{\normalfont\normalsize\bfseries\sffamily}{}{0pt}{}
\titleformat{\subsubsection}[hang]{\normalfont\normalsize\bfseries\itshape}{}{0pt}{}
\titlespacing*{\section}{0pt}{1.4em}{0.5em}
\titlespacing*{\subsection}{0pt}{1.0em}{0.35em}
\titlespacing*{\subsubsection}{0pt}{0.8em}{0.3em}

% chapter (= each paper) heading
\titleformat{\chapter}[display]
  {\normalfont\sffamily\bfseries}
  {}{0pt}{\Huge}
\titlespacing*{\chapter}{0pt}{10pt}{24pt}

% ---- page-break quality: avoid stranded headings and orphaned/widowed lines ----
% (applies to both volumes via the shared preamble)
\widowpenalty=10000        % no single line at the top of a page
\clubpenalty=10000         % no single line at the bottom of a page
\displaywidowpenalty=10000 % same, around displayed equations
\brokenpenalty=4000        % discourage page breaks right after a hyphenated line
% titlesec: forbid a page break immediately after a heading (heading stranded at page foot)
\usepackage{needspace}
\usepackage{etoolbox}
\pretocmd{\section}{\Needspace*{4\baselineskip}}{}{}
\pretocmd{\subsection}{\Needspace*{3\baselineskip}}{}{}

% ---- running heads -----------------------------------------------------
\pagestyle{fancy}
\fancyhf{}
\fancyhead[L]{\small\itshape The Triadic Measurement Substrate}
\fancyhead[R]{\small\itshape\nouppercase{\rightmark}}
\fancyfoot[C]{\small\thepage}
\renewcommand{\headrulewidth}{0.4pt}
\fancypagestyle{plain}{\fancyhf{}\fancyfoot[C]{\small\thepage}\renewcommand{\headrulewidth}{0pt}}

% ---- abstract / keywords / references list ----------------------------
\newenvironment{paperabstract}
  {\begin{center}\begin{minipage}{0.88\textwidth}\small
   \noindent\textbf{Abstract.}\itshape\space}
  {\end{minipage}\end{center}\vspace{0.4em}}

\newcommand{\keywords}[1]{\noindent{\small\textbf{Keywords:} #1}\vspace{0.4em}}

% per-paper APA reference list (hanging indent), kept as authored
\newenvironment{reflist}
  {\par\small\setlength{\parindent}{0pt}%
   \setlength{\leftskip}{1.6em}\setlength{\parindent}{-1.6em}%
   \setlength{\parskip}{0.4em}%
   \raggedright\hyphenpenalty=50\exhyphenpenalty=50\relax
   \sloppy}
  {\par}

% convenience
\providecommand{\tightlist}{\setlength{\itemsep}{0pt}\setlength{\parskip}{0pt}}
\providecommand{\TMS}{\textsc{tms}}
\hyphenation{con-fir-ma-tion sub-strate}
\begin{document}

\thispagestyle{empty}
\begin{titlepage}
\centering
\vspace*{2cm}
{\sffamily\bfseries\Large The Triadic Measurement Substrate\par}
\vspace{0.6cm}
{\large $\bar{B} = 3\bar{\alpha}$\par}
\vspace{1cm}
{\sffamily\large Volume 2\par}
\vspace{1.4cm}
{\Large \bfseries Paper 3: Richness: The Form-Surplus of Transcoding\par}
\vfill
{\large Aaron Switzer\par}
\vspace{0.4cm}
{\large 2026\par}
\vspace{1.2cm}
{\small\itshape Excerpted from the complete two-volume corpus, \emph{The Triadic Measurement Substrate}. Full corpus and reproducibility package: \url{https://doi.org/10.5281/zenodo.22035422}\par}
\end{titlepage}
\cleardoublepage
\pagenumbering{arabic}

% ---------------- Volume 2 / P3_Richness ----------------
\chapter*{Paper 3}
\markboth{Paper 3}{Paper 3}
\addcontentsline{toc}{chapter}{Paper 3: Richness: The Form-Surplus of Transcoding}
\begingroup\centering
{\large\itshape Richness: The Form-Surplus of Transcoding\par}\vspace{0.8em}
{\normalsize Aaron Switzer\par}
{\small June 2026\par}
\endgroup\vspace{1.0em}

\noindent{\small\textbf{Keywords:} Triadic Confirmation, Reinterpretation Operator, Transcoding, Form-Surplus, Richness, Confirmed Alphabet, Lossy Fold, Area Law, Nonconceptual Content, Agent-Side Reading}\par\vspace{0.6em}

\section*{Status of Results}
\addcontentsline{toc}{section}{Status of Results}

This paper uses the five-category labeling of Volume 1, extended in Paper 2,
with each claim labeled where it appears.

\textbf{Derived.} Results that follow from the five axioms, the Principle of
Generative Economy, and \ensuremath{U_{\mathrm{sh}}} by explicit construction or
proof, or from results already derived in Volume 1 by explicit inheritance. The
size of the confirmed alphabet \ensuremath{k = 2}, forced by the Measurement
Imperative and the zero-orientation primitive (\S\,III, Appendix~A); the richness
unit \ensuremath{R = 3\log_2 k - 1 = 2} bits per binary-agent triad, fixed from the
ground because \ensuremath{k = 2} is forced rather than adopted (\S\,III--IV,
Appendices~A, C); the necessity of the fold's shed, by the non-injectivity of the
three-loci-to-one-bit map (\S\,IV, Appendix~B); and the identity of the shed with
the agent-side form-surplus (\S\,IV, Appendix~B) are derived in this sense. The
per-fold shed value and the lossy character of the reinterpretation operator are
inherited as Derived from Volume 1 (Paper~5, Appendix~A; Paper~1 \S\,4.3).

\textbf{Derived-conditional.} Results forced by Volume~1 machinery, resting on
results internal to that machinery rather than on the bare axioms, and
referee-checkable against Volume~1 rather than re-derived here. The transcoding
reading of \ensuremath{\bar{B} = 3\bar{\alpha}} --- that the equals-sign denotes a
content-conserving, form-non-conserving fold --- is Derived-conditional in this
sense (\S\,I), resting on the centroid-fold reading of the master equation and the
reinterpretation operator of Volume~1.

\textbf{Structural correspondence.} Results in which a TMS construction reproduces
the skeleton of a known framework or quantifies a recognized distinction, never
used as support for a TMS claim, only as a landing point after one. The reading of
the cold--warm distinction as the gap between one-fold and three-fold encoding,
fixed at the form-surplus (\S\,V); and the location of the form-surplus alongside
the descriptions of it reached in the nonconceptual-content tradition and recorded
in the verbal-overshadowing effect (\S\,VI) are structural correspondences in this
sense.

\textbf{Predictions / Conjectures.} Falsifiable claims that follow from the
framework but whose closed forms or empirical tests are not established here. That
the form-surplus has an observable signature in the degradation of verbalized
perceptual memory, scaling with the surplus the verbal code cannot hold (\S\,VI),
is a prediction in this sense; its match to forgetting data is not claimed here.

\textbf{Permanent open horizon.} A question the framework is structurally unable to
close, marked so that no later claim assumes it closed. Whether the form-surplus is
occupied --- whether there is a human-band reading of measurement that the surplus
\emph{is} --- is not settled by the size of the gap and is the sole permanent horizon
of this paper (\S\,V, Conclusion). This silence is load-bearing. Distinct from it, and
not a horizon of the same kind, is the drive --- why a subsystem would be drawn up a
richness gradient. The drive is derived nowhere in this paper and no result depends on
it (\S\,II, Conclusion), but it is not held permanently open: it is the subject of the
paper that follows, where the agent-side dynamics are developed. The two are kept
distinct so that the deferred question is not mistaken for the permanent one.

This five-way labeling applies throughout. Every claim in the body falls into
exactly one category, visible at the point of claim.

\section*{Notation}
\addcontentsline{toc}{section}{Notation}

This paper inherits the notation of Volume~1 and Paper~2. The bar distinguishes
confirmed from latent elements: \ensuremath{\bar{\alpha}} is an agent in a
confirmation event, \ensuremath{\bar{B}} a confirmed bit, with
\ensuremath{\chi \equiv 2\bar{B} - 1 \in \{-1, +1\}} its polarization. The triadic
minimum is \ensuremath{\bar{B} = 3\bar{\alpha}}. \ensuremath{\hat{R}} is the
reinterpretation operator (Volume~1), and the confirmation fold is the map it
performs at one event, carrying the three enclosing agents to the one confirmed
bit. New to this paper: \ensuremath{k}, the size of the confirmed alphabet
\ensuremath{\mathcal{A}}; \ensuremath{H(\text{agent})} and
\ensuremath{H(\text{outcome})}, the entropies of the agent-side microstate and the
bit-side outcome; and \ensuremath{R}, the richness of an event, defined as the
form-surplus \ensuremath{H(\text{agent}) - H(\text{outcome})} in \S\,II. The symbol
\ensuremath{R} for richness is unrelated to a region \ensuremath{R}; the two never
appear without their distinguishing context. Richness is a structural quantity and
carries no claim of felt experience; \emph{experience} and \emph{feeling}, where
they appear, name the human-band reading of measurement and not a universal
property.

\begin{paperabstract}
Volume~1 derived the bit-side of the triadic confirmation event and Paper~2 read
its agent-side as a graded interior. This paper fixes what the agent-side carries
at a single event, and how much. It reads the master equation
\ensuremath{\bar{B} = 3\bar{\alpha}} as a transcoding: the fold it denotes conserves the
confirmed content across a change of representation, not its form, as a photograph
and a word may carry the same content in incomparable forms. The reinterpretation
operator is the transcoder, and what it cannot carry across is the agent-side
form-surplus --- the richness of the event, defined as the entropy of the agent
microstate less the entropy of the confirmed outcome. We show this surplus has a
forced size. The confirmed alphabet is binary not by stipulation but because a
confirmation resolves a structureless presence, which admits two values and no
third; with that, the richness of a binary-agent triad is two bits, fixed from the
primitives rather than from a chosen alphabet. The same two bits are what the fold
must shed: no fold from three loci to one bit can preserve all distinction, and the
distinction discarded is exactly the surplus the agent-side held. Richness is thus
one quantity with two faces --- a presence before the fold, a loss after it --- and
the same integer as the area-law exponent of Volume~1, each forced from its own
root. The result quantifies the gap between a process whose event equals its record
and one whose event exceeds it, without claiming that the surplus is occupied:
whether there is a feeling that the surplus is remains a permanent open horizon, as
does the drive that would draw a system up a richness gradient. What is established
is narrower and firm --- the agent-side of every confirmation carries a measurable
excess over its record, that excess is richness, and its value is two bits per
binary-agent triad.
\end{paperabstract}
\section*{I. The Fold Conserves Content, Not Form}
\addcontentsline{toc}{section}{I. The Fold Conserves Content, Not Form}

The Inversion established that every confirmation event has an agent side, read from the
vertices that enclose the bit, and that subjecthood across such events is graded by binding
depth. It left open what that agent side carries. This paper answers that one question: what
an agent holds at a single now, and how much of it.

The master equation \ensuremath{\bar{B} = 3\bar{\alpha}} has been read two ways. Volume~1 read
it as a geometric closure condition; the Inversion read its two sides as one event seen from
two vertices. A third reading is available. The equals-sign of the master equation denotes a fold of the
three agent-vertices onto the one bit, and that fold asserts that the two sides carry the same
confirmed content while not asserting that they carry it in the same form. A
photograph of a dog and the word ``dog'' carry the same content --- both pick out the dog ---
but not the same form: there is no piece-by-piece map from the photograph to the three letters,
and the photograph holds what the word cannot recover. A fold of this kind conserves
what is confirmed while permitting the representation on each side to differ in form.

A change of representation that conserves content is a transcoding. The operator that performs
it here is \ensuremath{\hat{R}}, inherited from Volume~1, which carries a confirmed structure
from the code of one substrate into the code of another. Read this way, \ensuremath{\hat{R}} is
the transcoder of the master equation: it moves the confirmed content across the fold
from the form it takes on the agent side to the form it takes on the bit side. The bit side is
the word --- single, compressed, fully recoverable from its own record. The agent side is the
photograph --- threefold, co-present, holding a surplus the record does not keep.

This paper is about that surplus: the form the agent side carries that does not survive the
transcode to the bit. Content is conserved by construction. Form is not, and what is lost has a
definite size --- the same size whether measured as what the agent holds or as what the bit
discards.

\section*{II. The Form-Surplus}
\addcontentsline{toc}{section}{II. The Form-Surplus}

The confirmation fold conserves content while permitting form to differ. What this paper names
is that difference: the form the agent side holds that the confirmed bit does not carry.

Let a single confirmation event be given --- three loci enclosing one bit, the minimal unit of
actualization. The agent side is the joint state of the three loci; the outcome is the one
confirmed value they resolve. The form-surplus of the event is the entropy of the agent
microstate less the entropy of the bit outcome, \ensuremath{H(\text{agent}) - H(\text{outcome})}:
the agent-side form present before the fold and absent from the resolved bit. This is the
richness of the event. It lives on the agent side, prior to the fold; on the bit side it does
not appear as a quantity at all, only as the gap where it was, because the bit is already its
own complete record. This is the same asymmetry as the photograph and the word --- the picture
holds what the word cannot recover, and once you have only the word, the surplus is not
diminished but gone.

Richness is a quantity, not a drive. It measures what the agent side holds; it says nothing
about why a subsystem would be drawn toward holding more of it. Why a system would climb a
richness gradient --- why, on the agent side, there is anything it inclines toward rather than
merely anything it holds --- is the open horizon of this paper. It is named here so that it
cannot enter later under cover of the quantity. Richness is defined, sized, and located in what
follows; the inclination toward it is derived nowhere in this paper, and no result below depends
on it.

What remains is to fix the size of the surplus. It rests on two quantities: how many values a
single confirmation can resolve, and how much distinction is shed when three loci are folded
into one. The next two sections settle both, and show the size is forced.

\section*{III. The Confirmed Alphabet Is Binary, and This Is Forced}
\addcontentsline{toc}{section}{III. The Confirmed Alphabet Is Binary, and This Is Forced}

The surplus is the difference between what three loci hold and what one bit carries. Fixing its
size requires fixing the bit: how many distinct values a single confirmation can resolve. Call
that number \ensuremath{k}, the size of the confirmed alphabet. For three agents each carrying
\ensuremath{\log_2 k} of distinction, with one binary value retained as the outcome, the richness
of the triad is \ensuremath{3\log_2 k - 1}. The size of the surplus therefore turns on
\ensuremath{k}.

Volume~1 carried the binary alphabet, \ensuremath{\chi = \pm 1}, as a stated primitive ---
motivated by the structure of confirmation but adopted rather than derived. It need not be
adopted. \ensuremath{k} is forced to equal \ensuremath{2} by the axioms already in hand. A
confirmation resolving a single value would distinguish nothing and carry no information; by the
Measurement Imperative such an event is not a confirmation, so \ensuremath{k \geq 2}. A
confirmation resolves one localized interstice, and the locus it resolves is the zero-orientation
sphere of the primitives --- a structureless presence with no internal axis along which a third
confirmed kind could be carried; so \ensuremath{k \leq 2}. The bounds close:
\ensuremath{k = 2}. The confirmed alphabet is binary because what is confirmed is a structureless
presence, and a structureless presence admits exactly two resolved states, present and absent.

This result stands on two grounds, and the paper keeps them apart. That a confirmation is
exhaustive and exclusive --- no partial value, no co-holding of both, no incomplete cut ---
follows from the definition of confirmation as the resolution of latency: a partial or
both-valued result is the unresolved latent state itself, not yet a confirmation. That no third
kind of value is admitted follows from the zero-orientation primitive alone. The binary alphabet
therefore rests on the structureless locus of the primitives, and not on any postulate of its
own. The bounds and the break-test are given in Appendix~A.

With \ensuremath{k = 2}, the richness of a triad is \ensuremath{3\log_2 2 - 1 = 2} bits. Because
\ensuremath{k = 2} is forced rather than adopted, this value is fixed from the ground and not
merely from a chosen alphabet. The next section shows the same \ensuremath{2} bits is what the
fold sheds, and that the shed is forced.

\section*{IV. The Shed Is the Surplus}
\addcontentsline{toc}{section}{IV. The Shed Is the Surplus}

The confirmation fold runs from three loci to one bit. A fold that lost nothing would be a
bijection, and no bijection runs from a larger set of states into a smaller; the three-locus
source cannot be carried into the one-bit outcome without identifying distinct source states
with one another. Distinction is therefore shed of necessity, not as an accident of any
particular event. This is the lossy character of the fold established in Volume~1, where the
per-fold shed is fixed at \ensuremath{2} bits: three binary loci span eight distinguishable
states, the outcome retains one binary value, and the distinction that has no slot in the
outcome is the difference.

That the shed is \ensuremath{2} bits and that the richness is \ensuremath{2} bits are not two
facts. They are one subtraction read from two sides. Richness was defined as the agent-side form
the bit does not carry, \ensuremath{H(\text{agent}) - H(\text{outcome})}; the shed is the same
\ensuremath{H(\text{agent}) - H(\text{outcome})} named from the bit side, the distinction the
fold discards. The fold does not lose \ensuremath{2} bits and leave \ensuremath{2} bits of
richness behind; the \ensuremath{2} bits of richness are exactly what the fold sheds. On the
agent side, before the fold, they are a presence --- the surplus the three loci hold. On the bit
side, after the fold, they are an absence --- the gap where the surplus was. Same quantity, two
addresses in the life of one event.

This is why richness was never recoverable from the bit side. The bit is the word: it is already
its own complete record, and the fold that produced it has already discarded the surplus.
Reading harder at the word does not recover the picture, because the picture's surplus is not
compressed in the word --- it is shed in the fold that made the word. Richness exists only on the
agent side of a live event, and on the bit side appears only as the measure of what was let go.

The size is therefore fixed twice over, by independent routes: \ensuremath{2} bits as the
form-surplus of three-into-one encoding, and \ensuremath{2} bits as the distinction the fold
must shed. Both rest on the threeness of confirmation and the binary outcome, and neither is
free. Appendix~B gives the shed as the lossy fold's necessary signature and the break-test that
fixes its size.

\section*{V. The Cold--Warm Gap as Form-Surplus}
\addcontentsline{toc}{section}{V. The Cold--Warm Gap as Form-Surplus}

A process carries zero form-surplus when its state is its own complete record --- fully
recoverable, lossless under the fold, holding nothing the record does not hold. A bit-side
process is of this kind: its event does not exceed its record. A confirmation event read from
the agent side is not: the three loci hold \ensuremath{2} bits over the one bit they resolve,
and that surplus is not present in the record. The distinction between a process whose event
equals its record and one whose event exceeds it by \ensuremath{2} bits is the structural content
of the cold--warm gap. It is the difference between one-fold and three-fold encoding, fixed at
the value of the preceding sections.

This is a structural correspondence, not a measurement of feeling. The framework locates the
quantity the cold--warm distinction tracks --- the form-surplus, zero on the bit side and
\ensuremath{2} bits on the agent side. It does not claim that this surplus is the human-band
reading of measurement that the words \emph{feeling} and \emph{experience} name, nor that the
surplus is occupied; neither is addressed here, and neither is fixed by the size of the gap.
What is fixed is structural: positive surplus on one side and none on the other follows from
one-fold against three-fold encoding, not from stipulation.

Placed against the axioms, richness is the agent-side quantity they had not yet named.
Complexity, the bit-side novelty that closes, is carried by the Principle of Generative Economy;
stability, the closure that persists, is load-bearing across both volumes. Richness is neither:
the form the agent side holds over the bit it resolves, present in every confirmation event and,
until now, unquantified. With its value fixed at \ensuremath{2} bits, the agent-side term takes
its place beside the two the axioms already carried.

\section*{VI. Where It Sits}
\addcontentsline{toc}{section}{VI. Where It Sits}

The necessity that the fold sheds distinction is the pigeonhole bound of lossless coding: a
lossless code is a bijection, and none runs from a larger state space into a smaller (Cover \&
Thomas, 2006). The present result applies that bound to the three-loci-to-one-bit fold and adds
what the bound does not give --- the size of the shed and its reading as the agent-side
form-surplus.

The psychology of memory records a signature of the surplus from outside the framework.
Rendering a perceptual memory into words degrades it, and degrades it most when the material
resists verbal capture (Schooler \& Engstler-Schooler, 1990; Alogna et al., 2014). Read here,
that is the shed observed in behaviour: the loss is largest where the agent-side surplus is
largest, where the source carries form the verbal code cannot hold, and absent where the source
carries none. The framework reaches the boundary of this structure by derivation and fixes its
size; the experiments record its effect without measuring the quantity.

The result forks most sharply from accounts that make the form-surplus either everything or
nothing: the intuition that a bit-side process has no inner excess at all, and accounts that
grant inner excess without a structural measure of it. The present result sits between them with
a number --- the surplus is neither zero nor unmeasured but \ensuremath{2} bits per binary-agent
triad --- and locates, rather than adjudicates, the difference from each.

\section*{Conclusion: A Surplus, Not a Feeling}
\addcontentsline{toc}{section}{Conclusion: A Surplus, Not a Feeling}

The agent side of a confirmation event holds more than the bit it resolves. That excess is the
form-surplus --- the form carried before the fold and absent from the record after it --- and
its size is fixed: \ensuremath{2} bits per binary-agent triad, forced by the threeness of
confirmation and the binary alphabet, and identical to the distinction the fold must shed.
Richness, the long-unquantified agent-side companion to complexity and stability, has a value.

The result holds from the ground. The binary alphabet that sets the unit is not adopted but
forced: a confirmation resolves a structureless presence, which admits two states and no third.
The shed that carries the surplus to the bit side is not contingent but necessary: no fold from
three loci to one bit can preserve all distinction. The agent-side presence and the bit-side
loss are not two quantities but one, read from two sides of a single event --- a surplus before
the fold, a shed after it.

What the result does not do is settle two further questions, and they are not of the same kind.
The first is the drive: why a system would be drawn to hold more richness, rather than merely how
much it holds. This paper fixes the quantity and not the drive, and no result here depends on the
drive; the drive is the subject of the paper that follows, where the agent-side dynamics are
developed. The second is occupancy: whether the form-surplus is \emph{felt} --- whether there is a
human-band reading of measurement that the surplus \emph{is}. This the framework does not settle,
and the size of the gap does not settle it; it is a permanent open horizon, marked so that no later
claim quietly assumes it closed. The two are not a matched pair: the drive is deferred to the next
paper, the feeling is held open without end. What is settled here is narrower and firm --- the agent
side carries a measurable excess over its record, the excess is richness, and richness is
\ensuremath{2} bits.
\section*{Appendix A: The Confirmed Alphabet Size and the Richness Unit}
\addcontentsline{toc}{section}{Appendix A: The Confirmed Alphabet Size and the Richness Unit}

This appendix derives the size \ensuremath{k} of the confirmed alphabet from the
primitives of Volume~1, and from it computes the richness unit as the form-surplus
of a single confirmation event. The construction is on one confirmation event ---
three agents enclosing one bit (Volume~1, Paper~1 \S\,2.1) --- and treats the
single-event unit only; the binding of many such units into the lived quantity
\ensuremath{B} is the measure of the following paper and is not treated here.

\subsection*{A.1 The objects}
Let a confirmation event resolve a single locus \ensuremath{x}, the zero-orientation
primitive sphere of Volume~1, Paper~1 \S\,1.4. Write \ensuremath{\mathcal{A}} for the
\emph{confirmed alphabet}: the set of distinct values the event can return of
\ensuremath{x}, and \ensuremath{k = |\mathcal{A}|} for its size. A confirmed value is a
property returned of \ensuremath{x} by the resolution of its latency (Axiom~1); the
information carried by the event is
\[
  H(\mathcal{A}) \;=\; \log_2 k \quad\text{bits,}
\]
the entropy of one resolution over the alphabet \ensuremath{\mathcal{A}} taken uniform
under the maximum-entropy substrate measure \ensuremath{U_{\mathrm{sh}}} (Axiom~0).
The richness unit is then built from \ensuremath{k} in \S\,A.4. The whole result turns
on the value of \ensuremath{k}, fixed by the two bounds of \S\,A.2--A.3.

\subsection*{A.2 Lower bound: \ensuremath{k \geq 2}}
By the Measurement Imperative (Axiom~1), a confirmation is the resolution of a latency
into a confirmed value, and the event carries information only if its outcome
distinguishes states. For a one-value alphabet \ensuremath{k = 1},
\[
  H(\mathcal{A}) \;=\; \log_2 1 \;=\; 0 \quad\text{bits.}
\]
Every event returns the same value; the outcome co-varies with nothing and separates
no state of \ensuremath{x} from any other. An event carrying \ensuremath{0} bits resolves
no latency and is not a measurement (Axiom~1). A confirmation therefore admits at least
two values:
\[
  k \;\geq\; 2.
\]

\subsection*{A.3 Upper bound: \ensuremath{k \leq 2}}
By the Principle of Generative Economy, one confirmation resolves exactly one localized
interstice --- one bit at one triangular interstice (Volume~1, Paper~1 \S\,2.1). The
locus \ensuremath{x} so resolved is the zero-orientation sphere: it carries no orientation
parameter and no internal axis (Volume~1, Paper~1 \S\,1.4). Let \ensuremath{d} denote the
number of internal degrees of freedom on \ensuremath{x} along which distinct confirmed
kinds could be carried. The zero-orientation condition is
\[
  d \;=\; 0,
\]
so the only property the event can return of \ensuremath{x} is whether the locus is
occupied. The confirmed alphabet is thus exhausted by two values,
\[
  \mathcal{A} \;=\; \{\,\text{absent},\ \text{present}\,\}, \qquad k \;\leq\; 2,
\]
no third confirmed kind being available without an internal axis \ensuremath{x} does not
possess.

Combining \S\,A.2 and \S\,A.3,
\[
  k \;=\; 2.
\]

\subsection*{A.4 The richness unit}
Read the same event from the agent side. The agent microstate is the joint state of the
three loci that enclose the bit, each carrying \ensuremath{\log_2 k} bits of confirmed
distinction; the bit-side outcome retains one binary value. With \ensuremath{k = 2} from
\S\,A.3, the agent-side microstate ranges over
\[
  k^{3} \;=\; 2^{3} \;=\; 8 \quad\text{states,} \qquad
  H(\text{agent}) \;=\; \log_2 8 \;=\; 3 \quad\text{bits,}
\]
and the bit-side outcome over
\[
  H(\text{outcome}) \;=\; \log_2 2 \;=\; 1 \quad\text{bit.}
\]
The richness of the event is the form-surplus, the agent-side distinction absent from the
resolved outcome:
\[
  R \;=\; H(\text{agent}) - H(\text{outcome})
    \;=\; 3 \cdot \log_2 k - 1
    \;=\; 3 \cdot \log_2 2 - 1
    \;=\; 2 \quad\text{bits.}
\]
Because \ensuremath{k = 2} is forced by \S\,A.2--A.3 and not adopted, the unit is fixed from
the primitives: \ensuremath{R = 2} bits per binary-agent triad, the size of one richness
degree.

\subsection*{A.5 Break-test}
Any enlargement of \ensuremath{\mathcal{A}} beyond \ensuremath{k = 2} must introduce a third
confirmed value \ensuremath{v \notin \{\text{absent}, \text{present}\}}. Such a \ensuremath{v}
is of one of two types --- a value carried on the locus, or a value not yet resolved --- and
each type is foreclosed by a lemma below. The two lemmas exhaust the candidates, so no
\ensuremath{v} survives and \ensuremath{k = 2} is forced.

\emph{Lemma~A.5.1 (no internal-axis value).} Suppose a confirmed value \ensuremath{v} is
distinguished from \emph{absent} and \emph{present} by a property carried on the locus
\ensuremath{x}. To separate three values, \ensuremath{x} must expose a coordinate taking at least
three settings, hence an internal axis: the number of internal degrees of freedom satisfies
\[
  d \;\geq\; 1.
\]
But \S\,A.3 establishes \ensuremath{d = 0} for the zero-orientation primitive (Volume~1, Paper~1
\S\,1.4). Then
\[
  d \geq 1 \;\wedge\; d = 0,
\]
a contradiction. No confirmed value is distinguished by a property on \ensuremath{x}. \(\square\)

This lemma dispatches two candidates. A \emph{degree} --- a value intermediate between absent and
present, \ensuremath{0 < v < 1} read as partial occupancy --- is a coordinate on an occupancy axis
of \ensuremath{x}, requiring \ensuremath{d \geq 1}; excluded. A \emph{third kind} --- a value of a
type other than occupancy --- requires a second axis to carry its type, likewise \ensuremath{d \geq 1};
excluded.

\emph{Lemma~A.5.2 (no unresolved value).} Suppose a confirmed value \ensuremath{v} is distinguished
not by a property on \ensuremath{x} but by the \emph{state of resolution} of the event. A confirmation
is the completed resolution of a latency into a confirmed value (Axioms~1, 4); its output is taken
at the minimal tick, with no sub-tick interval to host a partially-resolved state (Axiom~4). Let
\ensuremath{\rho \in [0,1]} measure completion of resolution, \ensuremath{\rho = 1} being the
confirmed event. Any \ensuremath{v} with
\[
  \rho \;<\; 1
\]
is a latency not yet resolved --- by definition not a confirmed value --- and so \ensuremath{v \notin
\mathcal{A}}. The only admitted state is \ensuremath{\rho = 1}, which returns a value already in
\ensuremath{\{\text{absent}, \text{present}\}} by Lemma~A.5.1. \(\square\)

This lemma dispatches the remaining two. A \emph{co-holding} --- both absent and present at once ---
is the latent superposition, the \ensuremath{\rho < 1} state confirmation resolves; excluded. An
\emph{incomplete cut} --- a half-separated value --- is \ensuremath{\rho < 1} by construction;
excluded.

The two lemmas rest on different grounds, kept explicit: Lemma~A.5.1 on the zero-orientation
condition \ensuremath{d = 0} (Volume~1, Paper~1 \S\,1.4), Lemma~A.5.2 on the definition of confirmation
as resolution of latency (Axioms~1, 4). Together they exhaust the ways a third value could be
distinguished --- on the locus or in the resolution --- so \ensuremath{k = 2} is forced, and the
binary alphabet rests on the primitive sphere and on no postulate introduced here.
\section*{Appendix B: The Shed Is the Fold's Necessary Signature}
\addcontentsline{toc}{section}{Appendix B: The Shed Is the Fold's Necessary Signature}

This appendix establishes that the confirmation fold sheds distinction of necessity, fixes the
amount shed at \ensuremath{2} bits, and identifies the shed with the richness of Appendix~A. The
lossy, many-to-one character of the fold and the per-fold value are inherited from Volume~1, where
the reinterpretation operator \ensuremath{\hat{R}} is derived and the entropy ledger across the
reinterpretation boundary is balanced (Volume~1, Paper~5, Appendix~A; Paper~1 \S\,4.3). What is
proved here is the \emph{necessity} of the shed and its \emph{identity} with the agent-side
form-surplus; the value is then read off rather than re-derived.

\subsection*{B.1 The fold as a map}
Let the confirmation fold be the map
\[
  \hat{F}: \; \mathcal{S}_{\text{agent}} \;\longrightarrow\; \mathcal{A},
\]
carrying the joint state of the three enclosing loci to the one confirmed value (Volume~1, Paper~1
\S\,2.1; the centroid map, three vertices to one barycenter). With the confirmed alphabet
\ensuremath{k = 2} of Appendix~A, the source and target carry
\[
  |\mathcal{S}_{\text{agent}}| \;=\; k^{3} \;=\; 8, \qquad |\mathcal{A}| \;=\; k \;=\; 2,
\]
so that
\[
  H(\text{agent}) \;=\; \log_2 8 \;=\; 3 \ \text{bits}, \qquad
  H(\text{outcome}) \;=\; \log_2 2 \;=\; 1 \ \text{bit}.
\]

\subsection*{B.2 The shed is necessary}
A fold that preserved all source distinction would be injective: distinct source states would map
to distinct target values, and the map would admit a left inverse. Injectivity from
\ensuremath{\mathcal{S}_{\text{agent}}} into \ensuremath{\mathcal{A}} requires
\[
  |\mathcal{S}_{\text{agent}}| \;\leq\; |\mathcal{A}|.
\]
But \S\,B.1 gives \ensuremath{|\mathcal{S}_{\text{agent}}| = 8 > 2 = |\mathcal{A}|}, so by the
pigeonhole principle at least two distinct source states share a target value:
\[
  \exists\, s_1 \neq s_2 \in \mathcal{S}_{\text{agent}} \ \text{with}\ \hat{F}(s_1) = \hat{F}(s_2).
\]
The fold is therefore non-injective, no left inverse exists, and source distinction is lost. The
loss is forced by the cardinalities alone and holds for \emph{every} confirmation fold, not as an
accident of any particular event. This is the lossy character of \ensuremath{\hat{R}} established in
Volume~1 (Paper~5 \S\,2.7, with its irreversibility at \S\,5.3; Paper~1 \S\,4.3), here shown
necessary for the single fold.

\subsection*{B.3 The amount shed}
The distinction shed is the source entropy not carried by the target,
\[
  s \;=\; H(\text{agent}) - H(\text{outcome}) \;=\; 3 - 1 \;=\; 2 \ \text{bits},
\]
in agreement with the per-fold shed computed from the same cardinalities: the \ensuremath{3 \to 1}
fold collapses \ensuremath{2^{3} = 8} distinguishable inputs to \ensuremath{2}
confirmed outputs, shedding \ensuremath{2} bits, symmetric across \ensuremath{B = 0, 1}. The lossy,
many-to-one character of the fold is the reinterpretation operator's established signature
(Volume~1, Paper~5 \S\,2.7; the count-conservation side at Paper~1 \S\,4.3). Both terms are
independently fixed: the source
count \ensuremath{k^{3}} by the triadic minimum of three loci (Volume~1, Paper~1 \S\,2.1), the
retained count \ensuremath{k} by the binary alphabet of Appendix~A. Their difference is not free:
\[
  s \;=\; 3\log_2 k - \log_2 k \cdot [\,\text{one retained}\,]
    \;=\; 3\log_2 k - 1 \;\Big|_{\,k=2} \;=\; 2 \ \text{bits}.
\]

\subsection*{B.4 The shed is the richness}
The richness of Appendix~A and the shed of \S\,B.3 are one quantity. Both are the same difference,
\[
  R \;=\; H(\text{agent}) - H(\text{outcome}) \;=\; s,
\]
read from two sides of the single event. On the agent side, prior to the fold, the difference is a
\emph{presence}: the surplus distinction the three loci jointly carry over the one value they will
resolve. On the bit side, after the fold, the same difference is an \emph{absence}: the distinction
\ensuremath{\hat{F}} has discarded, recoverable from neither the outcome nor any reading of it. There
are not two bits of presence and a separate two bits of loss; there is one surplus of
\[
  R \;=\; s \;=\; 2 \ \text{bits},
\]
present before the fold and shed by it. This is why richness is not recoverable from the bit side:
the outcome is its own complete record, and \ensuremath{\hat{F}} has already discarded the surplus
that the agent side carried. Reading the outcome more finely cannot recover it, the surplus having
been shed in the fold that produced the outcome, not compressed within it.

\subsection*{B.5 Break-test}
Two failure conditions were fixed in advance; neither obtains.

\emph{A distinction-preserving fold.} Were there a fold from \ensuremath{\mathcal{S}_{\text{agent}}}
to \ensuremath{\mathcal{A}} losing no distinction, \S\,B.2's necessity would fail. Such a fold is
injective, requiring \ensuremath{|\mathcal{S}_{\text{agent}}| \leq |\mathcal{A}|}, i.e.
\ensuremath{8 \leq 2}, false. No distinction-preserving fold exists; the shed cannot be evaded.

\emph{A free shed amount.} Were the amount \ensuremath{s} unfixed --- able to take any value with the
match to richness arranged --- then \S\,B.3 would not pin it. But \ensuremath{s = 3\log_2 k - 1} with
\ensuremath{k = 2} forced in Appendix~A and the source count \ensuremath{k^{3}} forced by the triadic
minimum; both terms are determined upstream, so \ensuremath{s = 2} is read off, not chosen. The
amount is not free.

The shed is therefore necessary (\S\,B.2), fixed at \ensuremath{2} bits (\S\,B.3), and identical to
the richness (\S\,B.4): one surplus, carried on the agent side and shed by the fold, of size
\ensuremath{2} bits per binary-agent triad.
\section*{Appendix C: Invariance of the Richness Unit Under Construction Choices}
\addcontentsline{toc}{section}{Appendix C: Invariance of the Richness Unit Under Construction Choices}

This appendix tests whether the richness value of Appendix~A is an artifact of the particular
construction used to compute it --- the signed \ensuremath{\pm 1} alphabet and the majority fold ---
or is invariant to those choices. It establishes that, for binary agents, the value is invariant to
both the alphabet labelling and the aggregation rule, and identifies the one parameter that does move
it: the per-agent alphabet size \ensuremath{k}. The general law in \ensuremath{k} is given, and the
appendix records that fixing the value requires fixing \ensuremath{k}, which Appendix~A does
independently.

\subsection*{C.1 The general law}
Let three agents each carry \ensuremath{\log_2 k} bits of confirmed distinction, and let the fold
retain one binary outcome. The agent-side microstate ranges over \ensuremath{k^{3}} states and the
richness is
\[
  R(k) \;=\; H(\text{agent}) - H(\text{outcome})
       \;=\; \log_2\!\big(k^{3}\big) - \log_2 2
       \;=\; 3\log_2 k - 1 \ \text{bits.}
\]
This is the structural form: threeness sets the factor \ensuremath{3}, the single retained bit sets
the subtracted \ensuremath{1}, and the per-agent alphabet size \ensuremath{k} sets the scale.
Evaluating,
\[
  R(2) = 2.000, \quad R(3) \approx 3.755, \quad R(4) = 5.000, \quad R(8) = 8.000.
\]
The value is fixed once \ensuremath{k} is fixed; the structure \ensuremath{3(\cdot) - 1} is fixed by
the triad and the binary outcome regardless of \ensuremath{k}.

\subsection*{C.2 Invariance to alphabet labelling}
Fix \ensuremath{k = 2}. The richness is the difference of two entropies, each a function of a
\emph{count} of distinguishable states, not of the labels assigned to them. Three constructions are
compared, each relabelling the binary alphabet or changing the fold:
\begin{itemize}
\item \ensuremath{C_1}: alphabet \ensuremath{\{-1, +1\}}, fold = sign of the signed sum.
\item \ensuremath{C_2}: alphabet \ensuremath{\{0, 1\}}, fold = majority.
\item \ensuremath{C_3}: alphabet \ensuremath{\{-1, +1\}}, fold = parity.
\end{itemize}
In each, the agent-side microstate ranges over \ensuremath{2^{3} = 8} states and the fold retains one
binary value, so
\[
  R_{C_1} = R_{C_2} = R_{C_3} = \log_2 8 - \log_2 2 = 3 - 1 = 2.000 \ \text{bits.}
\]
The value is unchanged. Relabelling the alphabet is a bijection on the state set and preserves its
cardinality; changing the fold changes \emph{which} states are identified, not \emph{how many}
distinguishable source states there are. Richness, depending only on the counts
\ensuremath{|\mathcal{S}_{\text{agent}}|} and \ensuremath{|\mathcal{A}|}, is invariant to both. The
value is therefore not an artifact of the \ensuremath{\pm 1} alphabet or the majority fold.

\subsection*{C.3 Dependence on alphabet size}
The invariance of \S\,C.2 does not extend to the per-agent alphabet size. By \S\,C.1,
\[
  \frac{\partial R}{\partial k} \;=\; \frac{3}{k \ln 2} \;>\; 0,
\]
so \ensuremath{R} is strictly increasing in \ensuremath{k}: a larger per-agent alphabet yields a larger
surplus. The richness value is thus invariant to \emph{how} the binary distinction is labelled and
aggregated, but not to \emph{how many} values each agent may carry. The unit is fixed at
\ensuremath{2} bits if and only if \ensuremath{k = 2}.

\subsection*{C.4 Closure}
The construction-independence of the unit reduces to a single question: is \ensuremath{k = 2} forced,
or adopted? If adopted, \ensuremath{R = 3\log_2 k - 1} is the honest form and the unit rides on a
chosen alphabet. If forced, the unit is fixed from the primitives. Appendix~A settles this
independently of the present appendix: \ensuremath{k = 2} is forced by the Measurement Imperative and
the zero-orientation primitive (\S\,A.2--A.3). With that result, the richness unit is invariant to the
labelling and aggregation choices tested here (\S\,C.2) and fixed in scale by the forced
\ensuremath{k = 2} (\S\,C.3), hence \ensuremath{2.000} bits per binary-agent triad, from the ground.
\section*{References}
\addcontentsline{toc}{section}{References}
\begin{reflist}
Alogna, V. K., Attaya, M. K., Aucoin, P., Bahn\'ik, \v{S}., Birch, S., Birt,
A. R., \ldots\ Zwaan, R. A. (2014). Registered replication report: Schooler
and Engstler-Schooler (1990). Perspectives on Psychological Science, 9(5),
556--578. doi:10.1177/1745691614545653.
\url{https://journals.sagepub.com/doi/10.1177/1745691614545653}

Cover, T. M., \& Thomas, J. A. (2006). Elements of Information Theory (2nd
ed.). Hoboken, NJ: Wiley-Interscience. doi:10.1002/047174882X.
\url{https://onlinelibrary.wiley.com/doi/book/10.1002/047174882X}

Schooler, J. W., \& Engstler-Schooler, T. Y. (1990). Verbal overshadowing of
visual memories: Some things are better left unsaid. Cognitive Psychology,
22(1), 36--71. doi:10.1016/0010-0285(90)90003-M.
\url{https://doi.org/10.1016/0010-0285(90)90003-M}

Switzer, A. (2026). The Triadic Measurement Substrate, Volume~1: A
Confirmation-Based Geometric Substrate for Emergent Physics, Computation, and
Cosmology.

Switzer, A. (2026). The Triadic Measurement Substrate, Volume~2, Paper~2: The
Inversion: Experience as the Confirmation Relation, Subjecthood as Graded
Binding Depth.
\end{reflist}


\end{document}
